\subsubsection{Bounded Operators, Extensions}\label{subsec:bounded-operators-extensions}

We now modify our definition of a bounded operator and the norm of a bounded operator to fit our more general definition of operator.

\begin{Definition}
\label{def:bounded-operator-specified-domain}
A linear operator $(\mathcal{D}(A),A)$ from $X$ to $Y$ is said to be \textbf{bounded} if there exists a constant $C$ such that
\begin{equation}\label{eq:bounded-domain-operator-norm-inequality}
\|A(x)\|_Y\leq C\|x\|_X\quad\text{for every }x\in\mathcal{D}(A).
\end{equation}
If no such $C$ exists, the operator is said to be \textbf{unbounded}.

The norm of a bounded operator is the smallest $C$ for which \eqref{eq:bounded-domain-operator-norm-inequality} holds; i.e., the non-negative number
\begin{equation}\label{eq:bounded-domain-operator-norm-definition}
\|A\|:=\sup_{\substack{x\in\mathcal{D}(A)\\\|x\|_X\ne0}}\dfrac{\|A(x)\|_Y}{\|x\|_X}=\sup_{\substack{x\in\mathcal{D}(A)\\\|x\|=1}}\|A(x)\|_Y.
\end{equation}

\end{Definition}

Note that if $\mathcal{D}(A)=X$, this definition agrees with Definition~\ref{def:bounded-linear-mapping}. Definitions of properties such as continuity extend readily to the new version of the definition.

The following lemma tells us why we need to consider unbounded operators only when we study infinite-dimensional spaces.

\begin{Lemma}
\label{lem:finite-dimensional-domain-operators-bounded}
All operators with finite-dimensional domain $\mathcal{D}(A)$ are bounded.
\end{Lemma}

\begingroup
\postdisplaypenalty=10000
\renewcommand{\endProof}{\par\nobreak\hfill$\blacksquare$\par\vspace{0.1cm}}
\begin{Proof}
Let $\{e_i\}_{i=1}^n$ be a basis for $\mathcal{D}(A)$. By Lemma~\ref{lem:linear-mapping-bounded-continuity-equivalence} it is enough to show that $A(x_j)\to0$ whenever $x_j\to0$. To see this note that we can write $x_j=\sum_{i=1}^n\alpha_{i,j}e_i$. Furthermore, whenever $\lim_{j\to\infty}x_j=0$, we have
\begin{equation}\label{eq:finite-dimensional-operator-coefficient-convergence}
\lim_{j\to\infty}\alpha_{i,j}=0.
\end{equation}
Thus,
\begin{equation}\label{eq:finite-dimensional-operator-continuity-proof}
\lim_{j\to\infty}A(x_j)=\lim_{j\to\infty}\sum_{i=1}^n\alpha_{i,j}A(e_i)=0.
\end{equation}

\end{Proof}
\endgroup

We now make explicit the idea of extending an operator.

\begin{Definition}
\label{def:linear-operator-extension}
An operator $(\mathcal{D}(\widetilde A),\widetilde A)$ is said to be an \textbf{extension} of $(\mathcal{D}(A),A)$ if $\mathcal{D}(A)\subseteq\mathcal{D}(\widetilde A)$ and
\begin{equation}\label{eq:linear-operator-extension-agreement}
A(x)=\widetilde A(x),\quad\text{for every }x\in\mathcal{D}(A).
\end{equation}

\end{Definition}

The following theorem states we can extend any bounded operator to all of $X$ without changing its norm.

\begin{Theorem}
\label{thm:bounded-operator-norm-preserving-extension}
Every bounded operator $(\mathcal{D}(A),A)$ can be extended to all of $X$ without changing its norm; i.e., there exists a bounded operator $(\mathcal{D}(\widetilde A),\widetilde A)$ that extends $A$ and such that $\mathcal{D}(\widetilde A)=X$ and $\|\widetilde A\|=\|A\|$.
\end{Theorem}

\begingroup
\postdisplaypenalty=10000
\renewcommand{\endProof}{\par\nobreak\hfill$\blacksquare$\par\vspace{0.1cm}}
\begin{Proof}
We will prove this in only two special cases, but these include most applications.

\textbf{Case 1.} \emph{The domain $\mathcal{D}(A)$ is dense in $X$.} In this case we simply define the operator on all of $X$ by continuity; i.e., for every $x\in X$ there exists a sequence $x_n\in\mathcal{D}(A)$ such that $x_n\to x$. Since our operator is bounded, the sequence $Ax_n$ is Cauchy, and hence has a limit. Thus, we can define\[Ax:=\lim_{n\to\infty}Ax_n.\]To see that this definition is unambiguous, note that for any other sequence $\widetilde x_n\in\mathcal{D}(A)$ that converged to $x$, we would have $x_n-\widetilde x_n\to0$. Thus, by continuity\[\lim_{n\to\infty}[Ax_n-A\widetilde x_n]=\lim_{n\to\infty}[A(x_n-\widetilde x_n)]=0.\]The extended operator has the same norm since\[x\mapsto\dfrac{\|Ax\|}{\|x\|}\]is continuous at $x\ne0$.

\textbf{Case 2.} \emph{$X=H$, a Hilbert space.} If $\mathcal{D}(A)$ is not closed, we extend the operator to $\overline{\mathcal{D}(A)}$ by continuity, as in case 1 above. Since $H$ is a Hilbert space we can now use the projection theorem and define the operator $(\mathcal{D}(\widetilde A),\widetilde A)$ with $\mathcal{D}(\widetilde A)=H$ by
\begin{equation}\label{eq:hilbert-domain-operator-projection-extension}
\widetilde A(x)=\begin{cases}A(x),&x\in\overline{\mathcal{D}(A)}\\0,&x\in\overline{\mathcal{D}(A)}^{\perp}\end{cases}
\end{equation}
and extend to the rest of $H$ by linearity. The proof that $\|\widetilde A\|=\|A\|$ is left to the reader.

When $X$ is a Banach space we have to do some fancy footwork to get around the lack of an orthogonal decomposition when $\overline{\mathcal{D}(A)}\ne X$. Since this rarely occurs in applications, we skip this case in the proof.
\end{Proof}
\endgroup

Because of Theorem~\ref{thm:bounded-operator-norm-preserving-extension}, the domain of a bounded operator is almost always assumed to be the entire space $X$.

In Chapter~\ref{chap:function-spaces}, we saw that $\mathcal{L}(X,Y)$, equipped with the operator norm, is a Banach space. In the following, we shall sometimes also be interested in pointwise convergence rather than norm convergence. We make the following definition.

\begin{Definition}
\label{def:operator-strong-convergence}
We say that $A_n\in\mathcal{L}(X,Y)$ \textbf{converges strongly} to $A\in\mathcal{L}(X,Y)$ if $A_nx\to Ax$ for every $x\in X$.
\end{Definition}

For example, let $X=Y=\ell^2$ and let $A_n$ be the operator which truncates a sequence after $n$ terms. Let $A_n$ converges strongly to the identity, but not in the sense of the operator norm.
